% 统一填写题目与参赛信息。
\gmcmsetup{
  title={广义标度律框架下大语言模型的\texorpdfstring{\newline}{}算力资源配置与能力前沿预测},
  year={2026},
  edition={第二十三届},
  problem={F}, % A--F；兼容保留，不印在附件3没有的封面栏位中。
  team={26105330291},
  school={中南大学},
  member-a={吴禄谦},
  member-b={宋柏辰},
  member-c={何沛宇}
}

% ---------------------------------------------------------------------------
% 图表题注格式：标签（“图 4.1”“表 4.1”）加粗，标签与题注之间用全角冒号。
% 放在 config.tex 而不是 gmcm2026.cls，是为了让模板类文件保持原样、
% 本队的排版偏好集中在本文件，便于与 GMCM2026-Overleaf 同步。
% ---------------------------------------------------------------------------
\DeclareCaptionLabelSeparator{gmcmcolon}{：}
\captionsetup{labelfont=bf, labelsep=gmcmcolon}
\captionsetup[subfigure]{labelfont=bf}
% 图题只保留主题；面板解释与结果解读放在引用该图的正文中。
\captionsetup[figure]{justification=centering,singlelinecheck=true,skip=6pt}
\setlength{\emergencystretch}{1em}
% Use the verified bundled heading font for contest titles.
\renewcommand{\gmcmcontestfont}{\heiti}
\setmonofont[Path=fonts/,Scale=0.90]{DejaVuSansMono.ttf}
% Generated code excerpts contain Unicode mathematical symbols in comments.
% Render them explicitly; Latin Modern Mono does not provide these glyphs.
\lstset{literate=
  {→}{{\ensuremath{\rightarrow}}}1
  {⇒}{{\ensuremath{\Rightarrow}}}1
  {∂}{{\ensuremath{\partial}}}1
  {∈}{{\ensuremath{\in}}}1
  {≈}{{\ensuremath{\approx}}}1
  {α}{{\ensuremath{\alpha}}}1
  {β}{{\ensuremath{\beta}}}1
  {γ}{{\ensuremath{\gamma}}}1
  {λ}{{\ensuremath{\lambda}}}1
  {μ}{{\ensuremath{\mu}}}1
  {η}{{\ensuremath{\eta}}}1
  {κ}{{\ensuremath{\kappa}}}1
  {τ}{{\ensuremath{\tau}}}1
  {Δ}{{\ensuremath{\Delta}}}1
}

% 私人信息可写入 config.local.tex，格式同上；该文件被 Git 忽略。
% 手工打包不受 Git 忽略规则保护，分享前请检查压缩包。
% 默认封面使用附件3的原始Logo，自动填写以上学校、队号及姓名。
% 若需直接使用Word原版，可填写 docs/附件3*.doc 后导出 A4 封面，
% 并在 config.local.tex 设置：
% \gmcmsetup{cover-pdf={docs/cover.pdf}}
% 此 PDF 的第 1 页会原样替换默认封面；其中的信息不由上述字段自动填入。
\InputIfFileExists{config.local.tex}{}{}
